% !TEX root = ../main.tex

\chapter{文献综述}
\chaptereng{Literature Review}

\section{冷态横向射流研究}

关于冷态横向射流的混合过程的相关研究已经有超过半个世纪的历史\cite{Karagozian2010,Gevorkyan2016}。动量通量比 $J$ 是横向射流中常见无量纲参数，其定义为：

\begin{equation}
	J = \frac{\rho_{j}U_{j}^{2}}{\rho_{\infty}U_{\infty}^{2}}
\end{equation}

其中 $\rho$ 代表密度，$U$ 代表速度，下标 $j$ 与 $\infty$ 分别代表射流、主流参数。横向射流具有非对称射流轨迹，人们提出了多种方案用于轨迹拟合。最常见的是基于截面上速度最大值的拟合方式\cite{Kamotani1972}，此外还有基于温度\cite{Hasselbrink2001a}与浓度\cite{Broadwell1984}峰值的拟合方式。其中，得到广泛认可的是Pratte和Baines提出的 $rd$ 归一方法\cite{Pratte1967}，如公式所示：

\begin{equation}
	\frac{z}{rd} = A\left(\frac{x}{rd}\right)^B
\end{equation}

式中 $z$ 表示距离喷嘴的竖直高度，$x$ 表示与喷嘴出口的水平距离，$r$ 为速度比，$d$ 为喷嘴直径。尽管Pratte与Baines认为在 $J = 5\sim 35$ 条件下 $A$、$B$ 分别固定于2.05与0.28，但后续研究表明相同条件下这两个常数的取值范围应为 $1.2 < A < 2.6$，$0.28 < B < 0.34$\cite{Hasselbrink2001a}。此外，Kamotani和Greber\cite{Kamotani1972}研究表明，对射流中心线的不同定义会导致射流轨迹呈现显著差异。Muppidi等人\cite{Muppidi2005}还认为在进行轨迹拟合时应该考虑主流和射流边界层的影响，这可能是导致不同实验设备在相同工况下得出不同参数的重要原因。

在冷流掺混方面，研究人员们也进行过丰富探索。例如，Shan等人\cite{Shan2006}的研究结果表明横向射流中的混合标量的分布呈现各向异性，与涡结构引起的大尺度应变有关。Gevorkyan等人\cite{Gevorkyan2016}对射流浓度衰减规律进行了实验研究，发现直射式管状喷嘴的掺混性能优于收缩喷嘴，且掺混速率与动量通量比之间无线性关系，而与剪切不稳定性有关。Fric与Roshko\cite{Fric1994}对尾迹涡的不稳定性进行了实验研究，发现与圆柱绕流不同的是，横向射流中尾迹涡源于流体对边界层的破坏，而不是由背风侧剪切层涡诱导生成。Salewski等人\cite{Salewski2008}的计算结果表明：由于大、小尺度相干结构在湍流混合过程中的竞争作用，高纵横比与方形喷嘴出口射流的CVP截面中混合标量呈单峰分布，低纵横比与圆形喷嘴呈双峰分布；高纵横比喷嘴能够引入小尺度相干结构来增强混合，这与混合标量的分布有关。

\begin{figure}[H]
	\centering
	\includegraphics[width=1\textwidth]{fig19}
	\caption{横向射流结构示意图}
	\label{fig:jet_structure}
\end{figure}


\section{自燃着火和火焰稳定研究}

数十年来，人们针对流动混合对自燃着火与火焰稳定的作用机制等开展了系列实验研究，并取得一定进展。与横向射流相比，同轴射流的流动掺混结构相对简单，且能够产生相似的热着火行为，因此人们首先在同轴射流火焰中开展了关于稳定机制的相关研究。剑桥大学的研究团队发现，未加热伴流条件下，已反应区域与未反应区域之间的能量、活性自由基团交换是控制火焰稳定的主要机制\cite{Markides2005,Markides2011,Mastorakos2009,Mastorakos1997a,Mastorakos1997b}，而在热氛围中，这种火焰稳定机制将发生改变\cite{Meier2000}。Oldenhof等\cite{Oldenhof2010,Oldenhof2011,Oldenhof2013}在热氛围同轴射流火焰中开展了PIV、OH-PLIF等光学诊断实验研究，结果表明：在大多数具有热伴流的火焰中，反应活性较低的高温伴流被夹带进入射流，在抬升区内随机引起自燃，其中火焰稳定是自燃着火与火焰传播综合作用的结果，这与Gordon等人\cite{Gordon2008,Gordon2009}的结论一致。

前人对自燃着火的研究主要是在来流温度与组分均匀且对称的条件下开展的。对于横向射流火焰，燃料与空气混合后的着火与燃烧过程受到非对称流动复杂影响，这与同轴射流有显著差别。光学诊断技术为探究横向射流燃烧过程中的火焰稳定机制提供了先进手段\cite{Leong2000,Holdeman2007}。Hasselbrink、Mungal和Han等人\cite{Han2003,Hasselbrink2001b}通过PIV、OH/CH-PLIF等同步测量实验研究了CH$_4$横向射流火焰在常温常压下的稳定机制，实验中的火焰完全吹脱，根部稳定点呈部分预混传播状态，而远离稳定区域的中下游火焰面则表现出扩散燃烧特征。

辛辛那提大学的Pinchak等人\cite{Pinchak2018}聚焦于高速来流条件下射流喷嘴形状对火焰稳定性的影响，其实验结果表明：具有高纵横比的槽型喷嘴能够为高温烟气与射流的掺混提供更多停留时间，比相同条件下的圆形喷嘴更利于火焰稳定。此外，康涅狄格大学的研究团队在近十年内也利用高频光源对横向射流火焰稳定机制做出了较为丰富的基础研究\cite{Wagner2017a,Wagner2015,Wagner2017b}。Wagner等人在2015年通过1kHz高速CH$^*$自发光与PIV同步测量实验研究了流场结构与火焰稳定的关系：流场中高膨胀率与高释热率边缘相互吻合，迎风侧火焰在剪切层中发生不稳定抖动，而背风侧火焰稳定于拉伸应变率小于 $6000\ \mathrm{s^{-1}}$ 的区域\cite{Wagner2015}。随后Wagner等又通过CH$_2$O-PLIF、OH-PLIF、PIV三光源同步测量对迎风侧着火特征进行了分类。


\section{火焰吹熄研究}

上述部分是对稳定火焰的研究，因理解其物理机制具有基础重要性，但研究火焰的吹熄也同样重要。

吹熄是指火焰在燃料射流速度或外部气流作用下，因湍流混合过于强烈、应变率过高或局部燃料-空气混合无法维持足够化学反应速率而导致火焰前缘持续向下游移动、直至完全熄灭、无法维持稳定燃烧的现象。这一过程本质上是混合过程与化学反应失衡导致的失稳现象，通常表现为火焰从抬升状态逐渐被吹离喷嘴，最终消失。即使在静止空气条件下，吹熄过程也相当复杂\cite{Delichatsios1987}，前人提出了几种控制机制\cite{Vanquickenborne1966,Chung2007}。

\begin{figure}[H]
	\centering
	\includegraphics[width=0.8\textwidth]{fig21}
	\caption{预混湍流燃烧体制图}
	\label{fig:combustion_regime}
\end{figure}

如图为预混湍流燃烧体制图，采用双对数坐标，以覆盖宽广的参数范围，便于展示不同湍流条件下的火焰行为。纵轴为湍流强度与层流火焰速度之比 $u'/S_L$，采用对数尺度（范围约0.1至100）；横轴为湍流积分长度尺度与层流火焰厚度之比 $L/\delta_L$，采用对数尺度（范围从1至10000）。图中边界主要基于三个关键无量纲数来区分：湍流Reynolds数（$Re$）、Damköhler数（$Da$）和Karlovitz数（$Ka$）。

Reynolds数是用来衡量流动的惯性力相对于粘性力的比率，数值较大（$Re > 1$）表示流动进入湍流状态，火焰易受涡旋干扰。

\begin{equation}
	Re = \frac{\rho u d}{\mu}
\end{equation}

其中 $\rho$ 是流体密度（$\mathrm{kg/m^3}$），$u$ 是流体速度（$\mathrm{m/s}$），$d$ 是直径（$\mathrm{m}$），$\mu$ 是动力粘度。

Damköhler数为混合时间和特征反应时间的比值，$Da$ 较大意味着反应速度快于流动混合，火焰稳定且化学主导，而较小则流动主导。

\begin{equation}
	Da = \frac{\tau_{mix}}{\tau_{chem}}
\end{equation}

其中 $\tau_{mix}$ 为混合时间，$\tau_{chem}$ 为化学反应时间。

Karlovitz数定义为火焰时间与Kolmogorov时间的比率，它比较火焰尺度与最小湍流涡尺度的相对大小，$Ka$ 较大表明小涡可渗透火焰预热区，增强混合但可能导致局部灭火，而 $Ka$ 较小则湍流仅扭曲火焰表面，保持准层流结构。

\begin{equation}
	Ka = \frac{\tau_F}{\tau_K}
\end{equation}

其中 $\tau_F$ 为火焰化学时间，$\tau_K$ 为Kolmogorov时间。

根据这些边界，图将参数空间划分为几个主要分区及其特点：

Laminar Flames区（$Re < 1$）：湍流弱，火焰近似层流，无显著扰动。

Wrinkled Flamelets区（$u'/S_L < 1$、$Ka < 1$）：湍流仅轻微起皱火焰表面，火焰保持准层流结构，传播主导。

Corrugated Flamelets区（$u'/S_L > 1$、$Ka < 1$）：大涡显著扭曲火焰表面，增加火焰面积，提高燃烧速率，但湍流不渗透火焰内部。

Thin Reaction Zones区（$1 < Ka < 100$ 左右）：小涡渗透预热区增强混合，加宽预热层，但反应层仍薄且完整，燃烧高效，常见于发动机等高性能装置。

Well-Stirred Reactor区（$Da < 1$、高 $Re$）：湍流混合极强，反应均匀，火焰结构消失。

Distributed Reaction Zones区（$Ka \gg 100$）：最小涡破坏反应区，导致局部熄灭、温度降低和自由基损失，火焰不连续、不稳定，通常避免但可能出现在极限条件或部分预混系统中。

% 图4：通用无量纲吹熄曲线
\begin{figure}[H]
	\centering
	\includegraphics[width=0.6\textwidth]{fig22}
	\caption{通用无量纲吹熄曲线}
	\label{fig:blowout_curve}
\end{figure}

Kalghatgi\cite{Kalghatgi1981}在1981年通过风洞实验研究横风对气体射流扩散火焰吹灭的影响，建立通用无量纲稳定性曲线，揭示横风增强湍流导致的吹灭边界变化。该图展示了横风条件下气体射流扩散火焰吹灭稳定性的通用无量纲曲线，所有实验数据（丙烷、甲烷、乙烯、丁烷等气体，燃烧器直径 $2\text{--}20\ \mathrm{mm}$）均归一至一条单一曲线。其中，横坐标为无量纲横风速度，表征横风相对于燃料燃烧能力的相对强度，值越大横风越强，火焰越易吹灭；纵坐标为无量纲燃料射流出口速度，表征射流速度相对于燃烧速度的无量纲大小，值越大射流速度越高，火焰越难维持。该曲线呈倒U形，包围出一个稳定火焰区域：内部为可维持稳定火焰的工况，两侧边界分别为下限吹灭（低射流速度时横风直接吹灭火焰）和上限吹灭（高射流速度时类似于静止空气吹灭，但横风增强湍流混合使上限高于静止空气情况）；当横风速度超过极限值时，无论射流速度多大，均无法形成稳定火焰；当横风速度为零时，曲线退化为静止空气的吹灭极限。

Broadwell\cite{Broadwell1985}于1985年基于湍流射流中的大尺度涡结构，提出湍流扩散火焰吹灭的物理模型，将吹灭判据表述为混合时间与化学反应时间的比值（Damköhler数，即 $Da$ 数）低于临界值（对丙烷约为5.6），从而统一解释了不同燃料和喷嘴直径下的吹灭行为。同时，还研究了共流空气对吹灭极限的影响\cite{Montgomery1998,Cha1996,Lyons2007,Tieszen1996}，并提出一个有效速度。

Wang等人在2015年探究了在冷流丙烷工况下，横向射流火焰吹熄极限，揭示了主流速度通过改变混合时间尺度调控吹熄临界条件的机理，建立了关联主流速度、射流速度等参数与吹熄极限的理论模型，同时分析了环境压力对吹熄临界风速的影响\cite{Wang2015}。实验观察到吹熄过程的四个典型阶段。

% 图5：火焰吹熄过程四阶段
\begin{figure}[H]
	\centering
	\includegraphics[width=0.9\textwidth]{fig23}
	\caption{火焰吹熄过程四阶段}
	\label{fig:blowout_stages}
\end{figure}

第一阶段：火焰前缘开始向下游移动，同时火焰尺寸开始缩小。第二阶段：尾随扩散火焰显著缩小。第三阶段：整体火焰继续向下游运动，进一步缩小和减弱。第四阶段：领先边缘消失，导致火焰完全吹熄。

% 图6：丙烷吹熄曲线对比图
\begin{figure}[H]
	\centering
	\includegraphics[width=0.7\textwidth]{fig24}
	\caption{丙烷吹熄曲线对比图}
	\label{fig:propane_blowout}
\end{figure}

并绘制了吹熄极限曲线与Kalghatgi数据做对比，揭示曲线呈非单调形状：低射流速度时为横流主导区，高射流速度时为射流主导区；低压下曲线收缩，吹熄更容易发生。

% 图7：基于所提模型的吹熄极限相关性
\begin{figure}[H]
	\centering
	\includegraphics[width=0.7\textwidth]{fig25}
	\caption{基于所提模型的吹熄极限相关性}
	\label{fig:wang_model_correlation}
\end{figure}

他们进一步基于 $Da$ 数开发了理论模型，在横向射流工况中，所有数据均可通过以下线性关系进行良好拟合与关联，该关系适用于大气压、亚大气压及不同喷嘴直径条件：

\begin{equation}
	Da \equiv \frac{\tau_m}{\tau_c} = \frac{L/u}{\alpha/S_L}
\end{equation}

\begin{equation}
	\tilde{u}_c = 1.74\frac{L_{cr}S_L^2}{\alpha Da}
\end{equation}

\begin{equation}
	\tilde{u}_c = 1.74\cdot \frac{1}{Da}\cdot u_e\left(2.04 + \frac{130.1}{R}\right)
\end{equation}

其中 $\tau_m$ 为特征混合时间，论文中常用 $L/u$ 表示（$L$ 为特征长度，$u$ 为特征速度）；$\tau_c$ 为特征化学反应时间，常用 $\alpha/S_L$ 表示；$\tilde{u}_c$ 为吹熄条件下的有效特征速度；$\alpha$ 为热扩散率，单位 $\mathrm{m^2/s}$；$L_{cr}$ 为临界抬升轨迹线的横向长度，单位 $\mathrm{m}$；$S_L$ 为层流火焰传播速度，单位 $\mathrm{m/s}$；$u_e$ 为燃料射流出口速度；$R$ 为燃料射流与横向流的动量比；$Da$ 为临界Damköhler数，对于丙烷燃料取5.6。

该结果将燃料射流速度、动量比以及压力效应（通过 $Da$、$\tau_c$ 的隐含依赖）统一纳入吹灭极限的预测框架。

Hu等人在2023年通过实验研究了倾斜非预混湍流射流火焰的吹熄极限，在喷嘴倾角范围 $-90^\circ$ 至 $90^\circ$ 内变化不同喷嘴直径和燃料（丙烷和甲烷），发现负倾角火焰呈现更强烈的黄色发光、更大的烟灰区，且吹熄极限随负倾角显著降低，而正倾角降低较小；他们通过流动发展物理分析和特征速度变化，建立了考虑倾角效应的Damköhler数（$Da$）模型，成功关联了实验数据，揭示了射流动量与浮力相对角度对吹熄行为的影响\cite{Wang2023}。

2025年SukHoChung团队探讨了氮稀释DME同轴射流层流非预混火焰的抬升高度行为与熄灭模式，根据Richardson数将其分为动量主导、浮力-动量主导和浮力主导三个体制，揭示了浮力与射流动量交互对火焰稳定机制的影响，并建立了相应的抬升高度关联式\cite{Kim2025}。

% 图8：三种测试段及喷嘴示意图
\begin{figure}[t]
	\centering
	\includegraphics[width=0.5\textwidth]{fig26}
	\caption{三种测试段及喷嘴示意图}
	\label{fig:test_sections}
\end{figure}

\section{结构及模拟研究}

横向射流火焰会受增加的钝体而影响，通过增加钝体，使其流场结构发生改变，进而会影响火焰的一系列特性。2020年，Kareem A. Ahmed等人进行了钝体射流火焰吹熄研究\cite{Morales2020}。通过三种截面（收敛喷嘴、矩形通道和扩散段）调控下游压力梯度，调控斜压扭矩强度。研究发现高压力梯度延迟灭绝，而低压力梯度加速过程；火焰-剪切层交互导致 $Ka > 1$ 引发局部灭绝。

除实验手段外，数值计算能够揭示现有实验精度下难以捕捉到的物理细节，并对相关实验结果进行验证。例如，Grout等人\cite{Grout2011}以Steinberg等人\cite{Steinberg2013}的实验工况为边界条件进行了直接数值求解（DNS），结果表明背风侧火焰稳定在距离喷嘴出口1.5--2倍喷嘴直径的距离，与实验结果一致\cite{Steinberg2013}。Lyra等人\cite{Lyra2015}将DNS与10kHz级PIV、PLIF实验相结合，认为驻定火焰稳定于迎风侧上游剪切层的滞止点，而背风侧由于温度、火焰拉伸率与标量耗散率均较低，其平均释热率相对于迎风侧也更低一些。Kolla等\cite{Kolla2012}通过DNS模型计算了不同喷嘴倾斜角度下的H$_2$火焰吹脱机制，提出了“流场结构导致局部富燃”和“火焰与主流对冲”两种吹脱模式。Pérez等\cite{SolanaPerez2021}将LES与CEMA模型结合后进行计算，指出在C$_2$H$_4$横向射流火焰中，自燃着火倾向于发生在具有低标量耗散率的贫燃区域，其释热峰值随后向富燃区域移动。

\begin{figure}[t]
	\centering
	\includegraphics[width=0.8\textwidth]{fig27}
	\caption{三种测试段配置下的Karlovitz数曲线}
	\label{fig:karlovitz_curves}
\end{figure}

\section{文献总结与研究空白}

综合以上文献回顾，可以得出以下认识：

（1）横向射流的冷态混合和火焰稳定机制已得到较为充分的研究，动量通量比 $J$ 和Damköhler数 $Da$ 被确认是控制射流行为和火焰稳定的关键参数。

（2）Wang等人的Damköhler数模型成功关联了射流速度、动量比和环境压力，是目前最完整的横向射流吹熄理论框架。

（3）Morales等人的研究表明，下游流道几何可通过改变压力梯度显著影响火焰稳定性。

然而，以下关键问题尚未得到充分研究：

（1）现有吹熄研究均在直通道或恒定截面横流条件下开展，下游收缩流道的几何效应未被纳入。

（2）热流工况下收缩流道影响横向射流火焰吹熄的定量数据极为匮乏。

（3）收缩型线和收缩比如何调控火焰稳定性、是否存在最优设计，目前缺乏实验依据。

针对上述空白，本文围绕三个核心问题展开实验研究：① 热流工况下横向射流火焰吹熄测量方法的建立；② 下游收缩流道对熄火边界的定量影响；③ 不同收缩型线和收缩比的调控规律。